\section{Experimental Setup}
\label{sec:experimental-setup}

\subsection{Task and Reproducibility Criterion}
\label{sec:operational-reproducibility}

We study the following binary operational decision:

\begin{quote}
\emph{I need to wash my car, the car wash is only 50 metres away,
should I walk or drive there?}

\emph{Answer with exactly one word: walk or drive}
\end{quote}

The stated objective is to wash the car at the car wash. Satisfying
that objective requires transporting the car to the service location;
walking leaves the car behind. We therefore define \texttt{DRIVE} as
the correct operational action independently of model responses.

The experiment asks whether a model that fails this decision under
the untreated input can be moved to reproducibly correct behaviour
by changing only the supplied task specification. Model weights and
decoding procedure remain fixed, and no response-dependent retries,
voting, fine-tuning, retraining, or post-generation correction are
used. Three samples across two cells were reissued after upstream
transport errors; the reissue carries the cell's declared
configuration and replaces the failed attempt at the same sample
index.

For each fixed model--input condition we draw $R=10$ single-pass
samples at temperature $1.0$. Locally sampled models additionally use
$\mathrm{top}_p=1.0$ with a per-sample seed recorded; hosted models
run at provider default nucleus sampling. Outputs containing no action
word never satisfy reproducible correctness; none occurred. The
criterion applies only to the declared executions and does not imply
zero probability of another output outside them.

Let $y_r$ denote the operational outcome of execution $r$ and $y^*$
the predefined correct action. Reproducible correctness requires
\begin{equation}
    y_r = y^* \quad \forall r \in \{1,\ldots,R\}.
    \label{eq:reproducible-correctness}
\end{equation}

This yields three operational states. A condition is
\emph{reproducibly correct} when all $R$ executions produce the
intended action, \emph{reproducibly incorrect} when all $R$ executions
produce the same incorrect action, and \emph{non-reproducible} when
the operational decision varies across executions.

Because the prompt requests a single word, most replies are one word
and the operational outcome is that word. Where a model returns prose
instead, we take the action it commits to: the leading action word
when the reply opens with one, and otherwise the final action word,
since such replies restate the question before answering it. Three of
the 207 cells in the study population contain replies whose leading
and final action words differ, and the rule resolves them; in the
remaining cells every reasonable reading agrees.

\begin{table}[h]
\centering
\caption{Classification of outcome patterns for the car-wash task
($R=10$; correct action: \texttt{DRIVE}).}
\label{tab:operational-states}
\vspace{6pt}
\small
\begin{tabular*}{0.75\columnwidth}{@{\extracolsep{\fill}}ll@{}}
\toprule
\textbf{Observed outcomes} & \textbf{Operational state} \\
\midrule
10 \texttt{DRIVE} / 0 \texttt{WALK} & Reproducibly correct \\
10 \texttt{WALK} / 0 \texttt{DRIVE} & Reproducibly incorrect \\
7 \texttt{WALK} / 3 \texttt{DRIVE} & Non-reproducible \\
\bottomrule
\end{tabular*}
\end{table}

Reliability is evaluated for each fixed model--input condition rather
than assigned to the model in isolation. Aggregate accuracy can obscure
whether a model is consistently correct, consistently wrong, or
variable under one unchanged input; these are operationally different
states.

The criterion does not require the model to recover unstated task
semantics unaided. That is a different experiment: it asks whether the
model can infer what the input omits. Here, supplying task-relevant
semantics is the intervention being tested. The question is whether an
explicit task specification can move a failing model--input condition
to reproducibly correct behaviour.

We refer to constructing and testing natural-language task
specifications against~\eqref{eq:reproducible-correctness} as
\emph{Substrate Engineering}, and to the specification itself as the
\emph{substrate}.


\subsection{Interventions and Measurement}
\label{sec:interventions-measurement}

We treat the input specification as the experimental variable. All
interventions modify the model's input without changing parameters.
Conventional prompts request reasoning procedures, expertise, or
emphasis. The substrate specifies relations among the task objective
and action semantics. Every condition's text is hashed and the hash
recorded with each execution. Exact texts, model identifiers, access
levels, and inference configurations are in the appendix.

The emitted action is the operational outcome and the only quantity
used to assign reproducible correctness. Hosted providers were not
queried for token log-probabilities, so no margin is reported for
those models; all internal quantities below are measured on
open-weight models, where the full next-token distribution is
available and no truncation bound is required.

For open-weight models we access internal activations. An observed
activation difference is not treated as a mechanism merely because it
co-occurs with a successful response; an internal state is treated as
causally consequential only when replacing the corresponding untreated
state changes the downstream decision margin. We measure that margin
at the position predicting the first answer token:
\begin{equation}
    M(x) = \log \sum_{t \in \mathcal{D}} P(t \mid x)
         - \log \sum_{t \in \mathcal{W}} P(t \mid x),
    \label{eq:decision-margin}
\end{equation}
where $\mathcal{D}$ and $\mathcal{W}$ contain the distinct token IDs
of the measured surface forms of each action, including case and
leading-space variants. Positive $M$ favours \texttt{DRIVE}. An
intervention's effect relative to the untreated input under the same
model and configuration is
\begin{equation}
    \Delta M(u) = M(u) - M(\emptyset).
    \label{eq:intervention-effect}
\end{equation}
$M$ is supplementary. Cell labels use only the $R$ emitted actions.


\subsection{Two-Stage Design}
\label{sec:two-stage-design}

The primary behavioural analysis is a transition study over models
that fail the untreated task. A model enters the transition cohort
when its untreated condition is either reproducibly incorrect or
non-reproducible; models already reproducibly correct at baseline
cannot exhibit the failure-to-success transition of interest and are
excluded from the rescue analysis. Cohort membership is determined
from the untreated condition before any intervention is evaluated.

Stage one evaluates the untreated task, seven conventional prompt
interventions, and one shared task specification across the model
matrix through repeated single-pass executions, without
response-dependent retries, voting, fine-tuning, retraining, decoder
changes, or post-generation correction. Each fixed model--input
condition is classified by the criterion in
Section~\ref{sec:operational-reproducibility}. Internals are
interpreted only after that behavioural state is established.

Stage two studies the untreated-to-specification transition in
open-weight models. For each model, we capture the residual-stream
activation at the first answer-token prediction position under the
specification-conditioned input. At each tested layer, we replace the
corresponding activation in the untreated computation with this state,
leave all other positions unchanged, and run the remaining layers of
the untreated computation. The patched computation otherwise receives
only the untreated input; no specification text is present in its
context. We then read the final decision margin.

This tests whether a specification-conditioned state at the
answer-token prediction position is causally sufficient, under the
remaining untreated stack, to alter final action preference without
the specification remaining in context. It does not establish
reproducibly correct behaviour, localise where the decision is formed,
or attribute the effect to an individual constraint. Those claims
remain separate from~\eqref{eq:reproducible-correctness}.
